\section{Thirteen integers using only the Jones ratio theorem}

There is a second, independent route to a small signed equation which does
not use Sun's 2026 reduction.  It combines the fixed-polynomial carry code
above with the literal Jones ratio construction and a new guard that
simultaneously certifies an interval and an accurate rational approximation.
The output has
\[
 \boxed{13\text{ unrestricted integers},\quad 1409\text{ binary operations},
 \quad\deg F_{13}=226334238192.}
\]
More precisely its degree is $370496t+84272$, where $t=610895$.
This is not a variable-count improvement over the companion ten-integer
construction.  It is an explicit smaller Jones-based alternative with no
2026 Lucas-reduction dependency.  The fourteen-integer variant has $1420$
operations and slightly smaller degree $226295141096$; both are retained.

\subsection*{Exact specification}
The full division-free recipe is
\path{signed_audit/generated/signed13_b4.json}; its source coefficients
are the same coherent $31$-variable quartic tuple as the companion
construction.  Set $a=2^{4269877}$ and use precisely the variables
\[
 (\epsilon,g,h,s,w,\phi,i,j,z_1,z_2,z_3,v,u)\in\mathbb Z^{13}.
\]
Put $b=2\epsilon^2$, $\omega=w^2$, and use the fixed-polynomial B4 front end
with $B=ab^4$, $t=610895$, $L=4t$, and $N=2B^{9t+2}$.  It produces $R$ and
\(c=bB^t\).  Retain the literal Jones ratio registers
\[
\begin{gathered}
 U=N^2\omega,\quad Y=N^2s,\quad M=RY,\quad P=2M^2U,\quad
 K=R+1+h(P-1),\quad A=M(U+1),\\
 B_1=2R+1,\quad W_1=b\omega,\quad C=B_1+W_1+\phi,
 \quad J=4A-5,\quad Q=3W_1C-2(W_1^2-1),\\
 D=(A^2-1)C^2+1,\quad E=iJDC^2,\quad F=(A^2-1)E^2+1,\\
 G=A+F(CD+1-A),\quad H=B_1+jC,\quad I=(G^2-1)H^2+1,\\
 a_0=(P^2-1)K^2+1,\qquad a_1=DFI.
\end{gathered}
\]
The new combined guard is
\[
 \boxed{\Delta=g(c-g)K^2-c^2(C-KY)^2>0.}
\]
Only $s\ge1$, $\phi\ge0$, and $\Delta>0$ remain to be encoded.  Set
\[
 d_1=3s-1,\quad d_2=3\phi+2,\quad d_3=3\Delta-1,
 \quad C_0=1,\quad C_r=C_{r-1}(1+4d_rC_{r-1}z_r^2),
 \quad C_*=C_3.
\]
Let $D_0=JF$, $T_0=QF+J(H-C)$, $Z=\epsilon wi z_1z_2z_3$, and
\[
\begin{gathered}
 S=ZD_0^5,\quad r=T_0+D_0v,\quad
 T=(r^3-D_0^2r-D_0^3)(r^2+23D_0^2),\\
 w_1=C_*,\quad w_2=C_*(P^2+K^2+2),\quad X=Su-T,\\
 A_0=S^2C_*,\quad A_1=S^2w_1^2a_0,\quad A_2=S^2w_2^2a_1,\\
 V=X^2+A_2-A_0-A_1,\quad W=V^2+4X^2A_2-4A_0A_1,\\
 \boxed{F_{13}=W^2-16V^2X^2A_2=0.}
\end{gathered}
\]
All displayed registers are polynomial abbreviations, not additional
unknowns.  The three-factor Pell stack, quintic intersective merger and
three-root norm have the full signed-domain proofs in
\path{signed_audit/SIGNED_B4_PROOF.md}; the new guard proof is in
\path{signed_audit/THIRTEEN_GUARD_PROOF.md}.

\subsection*{Why the guard is sound}
Since $c>0$, $\Delta>0$ forces $K\ne0$ and $0<g<c$.  The exact identity
\[
 c^2\bigl(K^2-4(C-KY)^2\bigr)
   =4\Delta+(c-2g)^2K^2
\]
gives the Jones approximation $4(C-KY)^2<K^2$.  Positivity of $C,Y$
then forces $K>0$; the square condition on $a_0$ excludes $h\le0$ by the
usual gap between consecutive squares.  Reflection of $i$ and a Pell
orbit modulo $CF$ restore positive $i,j$ without changing the conditions.
To invoke the unchanged \texttt{Nine.ratio\_sound} theorem literally, supply
its free ceiling parameter, in the proof only, as $G_{\rm aux}=8gK^2+1$.
The positive integer $K^2-4(C-KY)^2$ is at least one, so its original
margin is positive.  Thus the theorem gives the required power-of-two
base and central-binomial divisibility.  The recovered $g<c$ supplies
the carry-code bounds.

\subsection*{Why the canonical witness satisfies the stronger guard}
The exact source \texttt{Ratio.exists\_B25core} chooses $b=2^k$ and
\[
 \omega=2^{2R+1-k},\quad
 U=\frac{N^2}{b}2^{2R+1},\quad
 \rho=\frac{(U+1)^{2R}}{U^R}=Y+f.
\]
Expanding the binomial, the fractional tail is
\[
 f=\sum_{j=0}^{R-1}\binom{2R}{j}U^{j-R}
   <\frac{2^{2R}}U=\frac{b}{2N^2}<1.
\]
The other terms are the integer $Y$, and $N^2\mid Y$ as in the original
construction.  The literal \texttt{ratio\_core} estimate, with
$A=RY(U+1)$, yields
\[
 \left|\frac CK-\rho\right|
 \le\frac{2(Y+f)}{Y(U+1)}<\frac4U.
\]
Consequently, since $R\ge8$,
\[
 \left|\frac CK-Y\right|
 <\frac b{N^2}\left(\frac12+\frac4{2^{2R+1}}\right)
 <\frac b{N^2}.
\]
The direct witness code has $1\le g<c$, while $c,b\le N$.  Therefore
\[
 \frac{c^2(C-KY)^2}{K^2}<\frac{c^2b^2}{N^4}\le1
 \le g(c-g),
\]
which proves $\Delta>0$.  Neither $Y$ nor the base is changed by the later
product-Pell extension.  Choosing $k=1+2r$ sufficiently large makes
$b=2(2^r)^2$ and $\omega$ a positive square.  The three stack witnesses,
merger witness and norm root now complete the integer solution.

Thus ordinary mathematics formalizable in ZFC proves that $F_{13}$ has an
integer zero exactly when the source quartic does, hence exactly when
there is a finite ZFC refutation.  The accompanying source, B4, polynomial,
modular and exact-degree checks are diagnostics, not a new end-to-end
Lean certificate or a proof of consistency.
